\documentclass[uplatex,dvipdfmx]{jlreq}
\usepackage[fleqn,tbtags]{mathtools}
\usepackage{jlreq-deluxe}
\usepackage[noalphabet]{pxchfon}
\usepackage{stix2}
\usepackage{hira-stix}
\usepackage{graphicx}
\usepackage{booktabs}
\usepackage{longtable}
\usepackage{array}
\usepackage{url}
\usepackage{float}
\usepackage{tikz}
\usetikzlibrary{arrows.meta,positioning,shapes.geometric}

\setlength{\parindent}{1em}
\setlength{\parskip}{0.2em}

\newcommand{\code}[1]{\texttt{#1}}

\begin{document}

\title{担当箇所の基本設計・詳細設計\\
{\large 楽譜データ形式・4パート振り分け・発音情報仕様}}
\author{25G1053 齋藤翔太}
\date{2026年5月12日}
\maketitle

\section*{本書の目的}

本書は，チーム23「Arduino オーケストラ」における齋藤担当範囲の基本設計・詳細設計を
まとめたものである．対象タスクは表\ref{tab:wbs-scope}の4項目である．
作成済みの「楽譜データ・発音情報仕様書」で定めた内容をもとに，システム設計書として
機能，構成，インターフェース，状態遷移，ハードウェア条件，ソフトウェア構造，
処理フローが追える形に整理する．

\begin{table}[H]
\centering
\caption{本書が対象とするWBS}
\label{tab:wbs-scope}
\begin{tabular}{p{0.16\textwidth}p{0.76\textwidth}}
\toprule
\textbf{WBS} & \textbf{担当内容} \\
\midrule
116 & 楽譜データ形式 基本設計．4種類の楽器パートをArduinoのメモリに埋め込む方針を決める．\\
128 & 楽譜データ詳細設計．課題曲を4種類の楽器パートに振り分ける．\\
129 & 楽譜データ詳細設計．テンポ，音の高さ，拍をどんな形でデータに書くか決める．\\
130 & 楽譜データ詳細設計．ArduinoからPCへ送る音データの中身と送信タイミングを決める．\\
\bottomrule
\end{tabular}
\end{table}

\section{基本設計}

\subsection{設計方針}

本担当範囲では，楽器ノードが自分のパートの楽譜を保持し，指揮者ノードから届く
拍情報に合わせて発音情報をPCへ送る仕組みを設計する．
設計方針は次の通りである．

\begin{itemize}
  \item 楽譜はPC側ではなく，各楽器ノード（node\_02〜node\_05）のプログラムメモリに埋め込む．
  \item 楽譜はミリ秒ではなく，音高と拍を基準にした相対表現で持つ．
  \item テンポは楽譜に固定せず，指揮者ノードから届くBPMにより実時間へ変換する．
  \item 通信仕様は塩澤担当のArduino設計に合わせ，PCへ送るNOTEは20バイト固定長のバイナリ形式に統一する．
  \item 4つの楽器ノードは共通処理を使い，\code{partId}，\code{startBeatNo}，楽譜配列だけを差し替える．
\end{itemize}

\subsection{機能一覧とFBS対応}

齋藤担当範囲は，チーム全体FBSのうちF3「楽譜進行」とF4「発音情報配信」に対応する．
表\ref{tab:functions}に機能一覧を示す．

\begin{longtable}{p{0.12\textwidth}p{0.22\textwidth}p{0.52\textwidth}}
\caption{担当範囲の機能一覧}\label{tab:functions}\\
\toprule
\textbf{ID} & \textbf{機能} & \textbf{役割} \\
\midrule
\endfirsthead
\toprule
\textbf{ID} & \textbf{機能} & \textbf{役割} \\
\midrule
\endhead
F3.2 & 楽譜データ保持 &
各楽器ノードが自分のパートの楽譜をC配列として保持する．外部ストレージは使用しない．\\
F3.3 & BEAT同期での音符進行 &
指揮者ノードから届くBEAT番号とマスタ時刻に合わせ，発火すべき楽譜イベントを選ぶ．\\
F3.4 & NOTEイベント生成 &
音符開始時にNoteOn，音符終了時にNoteOffを生成する．休符ではNOTEを送信しない．\\
F3.5 & パート固有ロジック &
\code{partId}，\code{startBeatNo}，楽譜配列をノードごとに差し替え，4声の輪唱とリズムを構成する．\\
F4.1 & NOTEパケット生成 &
発音情報を\code{partId}，\code{noteNumber}，\code{velocity}，\code{gate}，
\code{durationMs}へ変換する．\\
F4.2 & Processingへの送信 &
20バイト固定長のNOTEをUSB SerialでPCへ送信する．\\
\bottomrule
\end{longtable}

\subsection{システム構成}

ユーザは指揮者ノードを振る．指揮者ノードはBPMとBEATを楽器ノードへ送る．
楽器ノードは自分の楽譜を進め，PC上のProcessingへ発音情報を送る．
Processingは受け取ったNOTEから音を合成し，スピーカへ出力する．

\begin{figure}[H]
\centering
\resizebox{\linewidth}{!}{%
\begin{tikzpicture}[
  node distance=10mm and 14mm,
  box/.style={draw, rounded corners, align=center, minimum width=31mm, minimum height=10mm},
  small/.style={draw, rounded corners, align=center, minimum width=25mm, minimum height=9mm},
  arrow/.style={-{Latex[length=2.2mm]}, thick}
]
\node[box] (user) {ユーザ\\指揮を振る};
\node[box, right=of user] (conductor) {node\_01\\指揮者ノード};
\node[box, right=of conductor] (wifi) {WiFi UDP\\CTRL / BEAT};
\node[small, above right=4mm and 13mm of wifi] (n2) {node\_02\\金管1};
\node[small, right=13mm of wifi] (n3) {node\_03\\金管2};
\node[small, below right=4mm and 13mm of wifi] (n4) {node\_04\\金管3};
\node[small, below=13mm of n4] (n5) {node\_05\\ドラム};
\node[box, right=23mm of n3] (pc) {PC / Processing\\音色合成};
\node[box, right=of pc] (spk) {スピーカ\\演奏音};

\draw[arrow] (user) -- (conductor);
\draw[arrow] (conductor) -- (wifi);
\draw[arrow] (wifi) -- (n2);
\draw[arrow] (wifi) -- (n3);
\draw[arrow] (wifi) -- (n4);
\draw[arrow] (wifi) -- (n5);
\draw[arrow] (n2) -- node[above, font=\small]{NOTE} (pc);
\draw[arrow] (n3) -- (pc);
\draw[arrow] (n4) -- (pc);
\draw[arrow] (n5) -- (pc);
\draw[arrow] (pc) -- (spk);
\end{tikzpicture}
}
\caption{ユーザ目線のシステム構成}
\label{fig:system}
\end{figure}

\subsection{操作インターフェースと状態遷移}

齋藤担当範囲では，専用の物理スイッチや画面は追加しない．
ユーザの操作は指揮者ノードを振ることに集約され，楽器ノードは通信状態と
楽譜進行状態に応じて動作する．

\begin{table}[H]
\centering
\caption{ユーザ操作と楽器ノードの挙動}
\label{tab:user-flow}
\begin{tabular}{p{0.24\textwidth}p{0.31\textwidth}p{0.36\textwidth}}
\toprule
\textbf{ユーザ操作} & \textbf{入力される情報} & \textbf{楽器ノードの挙動} \\
\midrule
電源を入れる & USB給電，WiFi接続開始 & \code{Idle}から\code{WaitSync}へ進み，CTRL受信を待つ．\\
指揮を振り始める & BEATとBPMが届く & \code{startBeatNo}に達したパートから演奏を開始する．\\
速く振る／遅く振る & CTRLのBPM値が変化する & 次の音符から新しいBPMで発音長を計算する．\\
BEATが一時的に欠ける & BEAT未受信時間がしきい値を超える & 直前BPMで\code{SelfRun}し，実BEAT復帰後は通常演奏へ戻る．\\
曲が終わる & 最終イベントまで発火 & デモ時は曲頭へ戻り，ループ演奏する．\\
\bottomrule
\end{tabular}
\end{table}

\begin{figure}[H]
\centering
\resizebox{0.98\linewidth}{!}{%
\begin{tikzpicture}[
  state/.style={draw, rounded corners, align=center, minimum width=27mm, minimum height=10mm},
  edgeLabel/.style={font=\scriptsize, fill=white, inner sep=1.2pt},
  arrow/.style={-{Latex[length=2mm]}, thick}
]
\node[state] (idle) at (0,0) {Idle\\起動直後};
\node[state] (sync) at (4.3,0) {WaitSync\\時計同期待ち};
\node[state] (wait) at (8.6,0) {WaitStart\\入り拍待ち};
\node[state] (play) at (13.2,0) {Playing\\演奏中};
\node[state] (finish) at (8.6,-2.8) {Finished\\曲末};
\node[state] (self) at (13.2,-2.8) {SelfRun\\自走中};

\draw[arrow] (idle) -- node[above, edgeLabel]{WiFi接続} (sync);
\draw[arrow] (sync) -- node[above, edgeLabel]{CTRL受信} (wait);
\draw[arrow] (wait) -- node[above, edgeLabel]{開始拍に到達} (play);
\draw[arrow] (play) to[bend left=18] node[right, edgeLabel]{欠損} (self);
\draw[arrow] (self) to[bend left=18] node[left, edgeLabel]{復帰} (play);
\draw[arrow] (play.south west) -- ++(-1.1,-0.9) -- node[above, edgeLabel]{曲末} (finish.east);
\draw[arrow] (finish) -- node[left, edgeLabel]{ループ} (wait);
\end{tikzpicture}
}
\caption{楽器ノードの状態遷移}
\label{fig:state}
\end{figure}

\subsection{4パート振り分け}

課題曲候補は\textbf{「かえるのうた」}とする．輪唱しやすく，テンポ追従の確認がしやすいためである．
4パートの割り当てを表\ref{tab:part-map}に示す．

\begin{table}[H]
\centering
\caption{4パート振り分け}
\label{tab:part-map}
\begin{tabular}{p{0.13\textwidth}p{0.15\textwidth}p{0.13\textwidth}p{0.12\textwidth}p{0.35\textwidth}}
\toprule
\textbf{ノード} & \textbf{partId} & \textbf{役割} & \textbf{入り拍} & \textbf{楽譜内容} \\
\midrule
node\_02 & \code{0x02} & 金管1 & 0拍 & 主旋律を曲頭から演奏する．\\
node\_03 & \code{0x03} & 金管2 & 4拍 & 主旋律を4拍遅れて演奏する．\\
node\_04 & \code{0x04} & 金管3 & 8拍 & 主旋律を8拍遅れて演奏する．\\
node\_05 & \code{0x05} & ドラム & 0拍 & 1拍ごとのリズム伴奏を演奏する．\\
\bottomrule
\end{tabular}
\end{table}

\section{詳細設計}

\subsection{ハードウェア設計}

本担当範囲で直接追加する部品はないが，楽譜データとNOTE送信が前提とする機器，
接続関係，動作条件を表\ref{tab:hardware}と図\ref{fig:hardware}に示す．

\begin{longtable}{p{0.20\textwidth}p{0.23\textwidth}p{0.47\textwidth}}
\caption{使用する機器と役割}\label{tab:hardware}\\
\toprule
\textbf{機器} & \textbf{型番・構成} & \textbf{役割} \\
\midrule
\endfirsthead
\toprule
\textbf{機器} & \textbf{型番・構成} & \textbf{役割} \\
\midrule
\endhead
楽器ノード & Arduino UNO R4 WiFi × 4 & \code{node\_02}〜\code{node\_05}．楽譜配列を保持し，NOTEを送信する．\\
指揮者ノード & Arduino系マイコン × 1 & CTRL / BEATをWiFi UDPで配信する．楽譜データは保持しない．\\
PC & MacまたはPC × 1〜4 & Processingを実行し，USB SerialからNOTEを受け取る．\\
接続ケーブル & USB Type-C等 & 楽器ノードとPCを1対1，またはUSBハブ経由で接続する．\\
音声出力 & PC内蔵または外部スピーカ & Processingが合成した音を出力する．\\
ストレージ & 追加なし & 楽譜はArduinoのプログラムメモリに埋め込むため，SDカードは使わない．\\
\bottomrule
\end{longtable}

\begin{figure}[H]
\centering
\resizebox{0.95\linewidth}{!}{%
\begin{tikzpicture}[
  node distance=10mm and 16mm,
  box/.style={draw, align=center, minimum width=34mm, minimum height=10mm},
  arrow/.style={-{Latex[length=2.2mm]}, thick}
]
\node[box] (n1) {node\_01\\CTRL / BEAT送信};
\node[box, right=of n1] (n) {node\_02〜05\\楽譜配列 + NOTE生成};
\node[box, right=of n] (pc) {PC\\Processing};
\node[box, right=of pc] (spk) {スピーカ};
\draw[arrow] (n1) -- node[above, font=\small]{WiFi UDP} (n);
\draw[arrow] (n) -- node[above, font=\small]{USB Serial 20B NOTE} (pc);
\draw[arrow] (pc) -- node[above, font=\small]{音声信号} (spk);
\end{tikzpicture}
}
\caption{ハードウェア接続関係}
\label{fig:hardware}
\end{figure}

\subsection{ソフトウェア構成}

楽器ノードは共通コードを使い，\code{ProjectConfig.h}と\code{score\_data.cpp}だけを
パートごとに差し替える．これにより，ノード間の処理差分を楽譜データと設定値に限定する．

\begin{verbatim}
firmware/
|-- node_02/
|   |-- include/
|   |   |-- ProjectConfig.h    // partId, startBeatNo, 楽器種別
|   |   `-- score_data.h       // ScoreEvent の宣言
|   `-- src/
|       |-- applyPattern.cpp   // 楽譜進行
|       `-- score_data.cpp     // node_02 用楽譜本体
|-- node_03/                   // node_02 と同構造
|-- node_04/
`-- node_05/
\end{verbatim}

\subsection{楽譜データ形式}

楽譜は\code{ScoreEvent}配列として定義する．小数の\code{float}拍ではなく，
\code{durationQ8}という固定小数表現を使う．\code{durationQ8 = 256}を1拍とし，
テンポ変更時はBPMから実時間へ変換する．

\begin{table}[H]
\centering
\caption{\code{ScoreEvent}のフィールド}
\label{tab:scoreevent}
\begin{tabular}{p{0.20\textwidth}p{0.20\textwidth}p{0.50\textwidth}}
\toprule
\textbf{フィールド} & \textbf{型} & \textbf{意味} \\
\midrule
\code{beatAt} & \code{uint16\_t} & このイベントを発火する拍番号．各パートの入り拍からの相対値．\\
\code{noteNumber} & \code{uint8\_t} & MIDIノート番号．\code{60}=C4．\code{0}は休符．\\
\code{velocity} & \code{uint8\_t} & 楽譜上の強さ．0〜127．CTRLのvelocityと合成する．\\
\code{durationQ8} & \code{uint16\_t} & 音の長さ．256が1拍，128が0.5拍，512が2拍．\\
\code{flags} & \code{uint8\_t} & bit0=NoteOn，bit1=NoteOff拡張，bit2=休符．\\
\bottomrule
\end{tabular}
\end{table}

\begin{verbatim}
struct ScoreEvent {
  uint16_t beatAt;      // 発火する拍番号（パート開始からの相対拍）
  uint8_t noteNumber;   // MIDIノート番号。0なら休符
  uint8_t velocity;     // 楽譜上の強さ 0〜127
  uint16_t durationQ8;  // 256 = 1拍
  uint8_t flags;        // bit0=NoteOn, bit2=休符
};
\end{verbatim}

\subsection{音高・拍・休符の表現}

\begin{table}[H]
\centering
\caption{楽譜データの表現ルール}
\label{tab:notation-rule}
\begin{tabular}{p{0.24\textwidth}p{0.66\textwidth}}
\toprule
\textbf{項目} & \textbf{ルール} \\
\midrule
音高 & MIDIノート番号で表す．C4は60，D4は62，E4は64，F4は65，G4は67とする．\\
休符 & \code{noteNumber = 0}かつ休符フラグを立てる．NOTEは送信せず，拍だけ進める．\\
拍 & \code{durationQ8}で表す．四分音符=256，八分音符=128，二分音符=512とする．\\
テンポ & 楽譜にはBPMを持たせない．CTRLで届くBPMから実時間へ変換する．\\
輪唱の入り & \code{ProjectConfig.h}の\code{startBeatNo}で指定する．\\
NoteOff & \code{durationQ8}から計算した時刻で自動生成する．\\
\bottomrule
\end{tabular}
\end{table}

\subsection{パート別設定}

\begin{verbatim}
struct ProjectConfig {
  uint8_t partId;          // 0x02〜0x05
  uint16_t startBeatNo;    // 輪唱の入り拍
  uint8_t instrumentKind;  // 1=金管, 2=ドラム
  bool loopEnabled;        // デモ時はtrue
};
\end{verbatim}

\begin{table}[H]
\centering
\caption{\code{ProjectConfig}のノード別値}
\label{tab:config-values}
\begin{tabular}{p{0.14\textwidth}p{0.16\textwidth}p{0.18\textwidth}p{0.18\textwidth}p{0.22\textwidth}}
\toprule
\textbf{ノード} & \textbf{partId} & \textbf{startBeatNo} & \textbf{instrumentKind} & \textbf{楽譜} \\
\midrule
node\_02 & \code{0x02} & 0 & 金管 & 主旋律A \\
node\_03 & \code{0x03} & 4 & 金管 & 主旋律A \\
node\_04 & \code{0x04} & 8 & 金管 & 主旋律A \\
node\_05 & \code{0x05} & 0 & ドラム & リズム伴奏 \\
\bottomrule
\end{tabular}
\end{table}

\subsection{楽譜データ例}

主旋律Aの冒頭は表\ref{tab:melody}のように表す．

\begin{table}[H]
\centering
\caption{主旋律A「かえるのうた」冒頭}
\label{tab:melody}
\begin{tabular}{cccccc}
\toprule
\textbf{順番} & \textbf{beatAt} & \textbf{音名} & \textbf{noteNumber} & \textbf{durationQ8} & \textbf{意味} \\
\midrule
1 & 0 & ド & 60 & 256 & 1拍鳴らす \\
2 & 1 & レ & 62 & 256 & 1拍鳴らす \\
3 & 2 & ミ & 64 & 256 & 1拍鳴らす \\
4 & 3 & ファ & 65 & 256 & 1拍鳴らす \\
5 & 4 & ミ & 64 & 256 & 1拍鳴らす \\
6 & 5 & レ & 62 & 256 & 1拍鳴らす \\
7 & 6 & ド & 60 & 256 & 1拍鳴らす \\
8 & 7 & 休符 & 0 & 256 & 1拍待つ \\
\bottomrule
\end{tabular}
\end{table}

\begin{verbatim}
const ScoreEvent kScore[] PROGMEM = {
  {0, 60, 96, 256, 0x01}, // ド
  {1, 62, 96, 256, 0x01}, // レ
  {2, 64, 96, 256, 0x01}, // ミ
  {3, 65, 96, 256, 0x01}, // ファ
  {4, 64, 96, 256, 0x01}, // ミ
  {5, 62, 96, 256, 0x01}, // レ
  {6, 60, 96, 256, 0x01}, // ド
  {7,  0,  0, 256, 0x04}  // 休符
};
\end{verbatim}

\subsection{ArduinoからPCへ送るNOTEデータ}

ArduinoからPCへ送るNOTEは，20バイト固定長のバイナリデータとする．
共通ヘッダ12バイトの後ろに，表\ref{tab:note-payload}のNOTEペイロード8バイトを置く．

\begin{table}[H]
\centering
\caption{NOTEペイロード}
\label{tab:note-payload}
\begin{tabular}{p{0.12\textwidth}p{0.20\textwidth}p{0.18\textwidth}p{0.40\textwidth}}
\toprule
\textbf{offset} & \textbf{field} & \textbf{type} & \textbf{意味} \\
\midrule
12 & \code{partId} & \code{uint8\_t} & \code{0x02}〜\code{0x05}．発音元パート．\\
13 & \code{noteNumber} & \code{uint8\_t} & MIDIノート番号．休符は送らない．\\
14 & \code{velocity} & \code{uint8\_t} & 0〜127．楽譜値とCTRL値の合成後．\\
15 & \code{gate} & \code{uint8\_t} & \code{1=NoteOn}，\code{0=NoteOff}．\\
16 & \code{durationMs} & \code{uint16\_t} & 発音予定長さ．Processing側の参考値．\\
18 & 予約 & \code{uint8\_t[2]} & 0で埋める．\\
\bottomrule
\end{tabular}
\end{table}

音符イベント発火時は\code{gate=1}のNoteOnを送る．
\code{durationQ8}ぶん経過した時点で\code{gate=0}のNoteOffを送る．
休符イベントではNOTEを送らず，楽譜ポインタだけを進める．

\subsection{内部API}

\begin{longtable}{p{0.31\textwidth}p{0.57\textwidth}}
\caption{楽譜進行で使用する内部API}\label{tab:api}\\
\toprule
\textbf{関数} & \textbf{役割} \\
\midrule
\endfirsthead
\toprule
\textbf{関数} & \textbf{役割} \\
\midrule
\endhead
\code{loadScoreEvent(index)} &
\code{PROGMEM}上の\code{kScore[index]}をRAM上の一時変数へ読む．\\
\code{isRest(event)} &
\code{noteNumber == 0}または休符フラグから休符か判定する．\\
\code{durationQ8ToMs(durationQ8, bpm)} &
\code{durationQ8}を現在BPMに基づきミリ秒へ変換する．\\
\code{mergeVelocity(score, ctrl)} &
楽譜上のvelocityと指揮者側CTRLのvelocityを合成する．\\
\code{emitNoteOn(event)} &
\code{pendingOn}を立て，NOTE送信モジュールへ渡すデータを作る．\\
\code{emitNoteOff()} &
発音中の音を止めるため，\code{pendingOff}を立てる．\\
\code{advanceScore(beatNo)} &
BEAT番号から発火対象イベントを順に処理する．\\
\code{buildNotePacket()} &
共通ヘッダとNOTEペイロードを20バイト列にシリアライズする．\\
\bottomrule
\end{longtable}

\subsection{処理フロー}

楽器ノードは入力，ロジック，出力の3フェーズで動作する．
楽譜データに関係する処理は\code{advanceScore()}に集約する．

\begin{figure}[H]
\centering
\resizebox{0.95\linewidth}{!}{%
\begin{tikzpicture}[
  node distance=8mm,
  proc/.style={draw, rounded corners, align=center, minimum width=53mm, minimum height=8mm},
  decision/.style={draw, diamond, aspect=2.3, align=center, inner sep=1mm},
  arrow/.style={-{Latex[length=2mm]}, thick}
]
\node[proc] (start) {起動・WiFi接続・Serial開始};
\node[proc, below=of start] (sync) {CTRL受信で時計同期};
\node[decision, below=of sync] (wait) {beatNo >=\\startBeatNo?};
\node[proc, below=of wait] (beat) {BEATをマスタ時刻で発火};
\node[decision, below=of beat] (event) {発火対象\\ScoreEventあり?};
\node[decision, below=of event] (rest) {休符?};
\node[proc, below left=8mm and 19mm of rest] (note) {NoteOn生成\\NoteOff時刻予約};
\node[proc, below right=8mm and 19mm of rest] (skip) {NOTE送信なし\\拍だけ進行};
\node[proc, below=22mm of rest] (off) {時刻到達でNoteOff生成};
\node[proc, below=of off] (send) {20バイトNOTEをSerial.write};
\node[proc, below=of send] (loop) {次BEATへ};

\draw[arrow] (start) -- (sync);
\draw[arrow] (sync) -- (wait);
\draw[arrow] (wait) -- node[right, font=\small]{はい} (beat);
\draw[arrow] (wait.east) -- ++(18mm,0) node[right, font=\small]{いいえ: 待機} |- (sync.east);
\draw[arrow] (beat) -- (event);
\draw[arrow] (event) -- node[right, font=\small]{はい} (rest);
\draw[arrow] (event.west) -- ++(-16mm,0) node[left, font=\small]{なし} |- (loop.west);
\draw[arrow] (rest) -- node[left, font=\small]{いいえ} (note);
\draw[arrow] (rest) -- node[right, font=\small]{はい} (skip);
\draw[arrow] (note) |- (off);
\draw[arrow] (skip) |- (loop);
\draw[arrow] (off) -- (send);
\draw[arrow] (send) -- (loop);
\draw[arrow] (loop.west) -- ++(-30mm,0) |- (beat.west);
\end{tikzpicture}
}
\caption{楽譜進行とNOTE送信の処理フロー}
\label{fig:flow}
\end{figure}

\subsection{BEAT受信時の処理}

\begin{enumerate}
  \item 受信したBEATの\code{beatNo}と\code{playAtMasterMs}をキューに入れる．
  \item \code{beatNo < startBeatNo}の場合は，自パートの入り前なので演奏しない．
  \item \code{effectiveBeatNo = beatNo - startBeatNo}を計算する．
  \item \code{kScore[currentEventIndex].beatAt <= effectiveBeatNo}のイベントを順に読む．
  \item 休符ならNOTEを送らず，\code{currentEventIndex}だけ進める．
  \item 音符なら現在BPMから発音長を求め，NoteOnを送る．
  \item 発音中の音について，予約した\code{noteOffAtMs}を過ぎたらNoteOffを送る．
\end{enumerate}

\subsection{テンポ変更時の処理}

楽譜データは拍単位で記録するため，テンポ変更時に配列を書き換える必要はない．
CTRLで新しいBPMを受け取った後，次に発火する音符から式\ref{eq:duration-ms}で実時間を計算する．

\begin{equation}
\label{eq:duration-ms}
durationMs = \frac{60000}{BPM} \times \frac{durationQ8}{256}
\end{equation}

すでに鳴っている音については，発音開始時に予約した\code{noteOffAtMs}を優先する．
これにより，途中でBPMが変わっても鳴りかけの音が不自然に切れない．

\subsection{例外処理}

\begin{longtable}{p{0.28\textwidth}p{0.28\textwidth}p{0.34\textwidth}}
\caption{例外・分岐時の処理}\label{tab:exceptions}\\
\toprule
\textbf{状況} & \textbf{判定条件} & \textbf{処理} \\
\midrule
\endfirsthead
\toprule
\textbf{状況} & \textbf{判定条件} & \textbf{処理} \\
\midrule
\endhead
入り拍前 & \code{beatNo < startBeatNo} & 演奏せず待機する．\\
休符 & \code{noteNumber == 0} & NOTEを送らず，次イベントへ進む．\\
BEAT欠損 & \code{beatTimeoutMs}超過 & 直前BPMから仮想BEATを作り\code{SelfRun}する．\\
遅延BEAT & 発音予定時刻を許容値以上過ぎた場合 & 発音をスキップし，ログに警告を出す．\\
同じBEATの重複 & \code{seq}または\code{beatNo}が処理済み & 二重発音を防ぐため捨てる．\\
曲末 & \code{currentEventIndex >= kScoreLength} & デモでは曲頭へ戻る．提出仕様では停止へ変更可能．\\
Serial送信失敗 & \code{Serial.write()}の戻り値が20未満 & 送信失敗ログを出し，次NOTEで復帰する．\\
\bottomrule
\end{longtable}

\section{関連資料との整合性}

\begin{longtable}{p{0.29\textwidth}p{0.31\textwidth}p{0.30\textwidth}}
\caption{関連資料との整合確認}\label{tab:consistency}\\
\toprule
\textbf{資料} & \textbf{記載内容} & \textbf{本書の対応} \\
\midrule
\endfirsthead
\toprule
\textbf{資料} & \textbf{記載内容} & \textbf{本書の対応} \\
\midrule
\endhead
\code{docs/design/score\_format.md} &
楽譜は各楽器マイコン内部で保持する． &
\code{ScoreEvent}配列をArduinoのプログラムメモリに埋め込む．\\
\code{docs/design/protocol.md} &
指揮者から楽器はUDP，楽器からPCは発音情報を送る． &
CTRL / BEATとNOTEを分け，NOTEはUSB Serial 20バイトで定義する．\\
ADR-0004 &
指揮1台，楽器4台，輪唱曲，金管楽器を想定． &
金管3声 + ドラム1声で「主旋律3 + リズム1」を満たす．\\
ADR-0007 &
指揮テンポ変更に演奏速度が追従する． &
楽譜を拍単位で持ち，BPMから実時間へ変換する．\\
\code{arduino\_塩澤.pdf} &
\code{ScoreEvent}，\code{startBeatNo}，20バイトNOTEを定義． &
同じ構造を採用し，結合時の形式不一致を避ける．\\
\bottomrule
\end{longtable}

\section{生成AIの利用について}

本書の作成には生成AIを利用した．主な利用目的は，作成済み仕様書の内容を基本設計・詳細設計
として再構成すること，ルーブリック項目への対応漏れを確認すること，および既存チーム資料との
整合性を確認することである．最終的な仕様判断は，チーム内資料とWBSに基づいて行った．

\end{document}
